\noindent
In particular, this lemma implies that a distinguished triangle
$(X, Y, Z, f, g, h)$ in $K(\mathcal{A})$ gives rise to a long exact
cohomology sequence
\begin{equation}
\label{equation-long-exact-cohomology-sequence-D}
\xymatrix{
\ldots \ar[r] &
H^i(X) \ar[r]^{H^i(f)} &
H^i(Y) \ar[r]^{H^i(g)} &
H^i(Z) \ar[r]^{H^i(h)} &
H^{i + 1}(X) \ar[r] & \ldots
}
\end{equation}
see (\ref{equation-long-exact-cohomology-sequence}). Moreover, there is
a compatibility with the long exact sequence of cohomology associated to
a short exact sequence of complexes (insert future reference here). For
example, if $(A^\bullet, B^\bullet, C^\bullet, \alpha, \beta, \delta)$
is the distinguished triangle associated to a termwise split exact
sequence of complexes (see
Definition \ref{definition-split-ses}),
then the cohomology sequence above agrees with the one defined using the
snake lemma, see
Homology, Lemma \ref{homology-lemma-long-exact-sequence-cochain}
and for agreement of sequences, see
Homology, Lemma \ref{homology-lemma-ses-termwise-split-long-cochain}.

\medskip\noindent
Recall that a complex $K^\bullet$ is {\it acyclic} if $H^i(K^\bullet) = 0$
for all $i \in \mathbf{Z}$. Moreover, recall that a morphism of complexes
$f : K^\bullet \to L^\bullet$ is a {\it quasi-isomorphism} if and only if
$H^i(f)$ is an isomorphism for all $i$. See
Homology, Definition \ref{homology-definition-quasi-isomorphism-cochain}.

